% GENERATED by `ggen sync` from `doc-pack` ontology — chapter/section/
% paragraph structure and prose live as RDF individuals, not hand-authored
% .tex. The ontology source is whichever file the invoking project's
% `ggen.toml` names under `[ontology].source` -- see this document's own
% `regenerate.sh`, which symlinks the actual `rdf/thesis.ttl` for this
% project into place. Do not edit this generated file by hand; edit the
% .ttl instance data and re-sync.
\chapter{Introduction}

\section{Thesis statement}
\label{sec:thesis-statement}


This paper consolidates two source theses into one account: \emph{Mission
Geometry: Planning, Conformance, and Separability for Bounded Receipted
Execution} (\texttt{docs/thesis/03\_planning\_geometry.tex}), and \emph{The
Projection Principle: A First-Principles Theory of Bounded Receipted
Execution for a Trillion Agents Beyond One Another's Comprehension}
(\texttt{docs/thesis/projection\_thesis.tex}). Read together, the two make
one claim: the manufacturing morphism $\Act = \muop(\Adm)$ and its chained
receipt are a single formal keystone, and PDDL planning, Petri-net
conformance, POWL separability, and lexicographic cost arbitration are its
concrete geometric instantiation, not a separate theory bolted alongside
it.


Every claim below traces to one of two artifact classes, cited by exact
path at the point of use: (a) the corpus RDF at
\texttt{docs/thesis/rdf/03\_planning\_geometry.ttl} and
\texttt{docs/thesis/rdf/projection\_thesis.ttl}, whose \texttt{math:statement}
literals are the source of every definition and theorem restated here; and
(b) the Lean~4 kernel receipts at
\texttt{tools/paper-factory/lean-pilot/formalization\_receipts.jsonl} (bare
Lean~4 core, no Mathlib) and
\texttt{tools/paper-factory/lean-lake/mathlib\_migration\_receipts.jsonl} (the
Mathlib-linked lane), whose \texttt{status} field is reported here exactly as
it reads on disk, in both directions: neither upgraded past what the kernel
accepted, nor downgraded past what it did.


\texttt{doc:register} marks every paragraph in this document
\textsc{evidence-bound} or \textsc{speculative}. \textsc{evidence-bound}
means every quantitative or factual claim in the paragraph names the
receipt line, corpus label, or source file that grounds it. Anything
motivated by, but not derived from or checked against, an artifact is
\textsc{speculative} and is quarantined into Chapter~\ref{ch:speculative},
never mixed silently into a technical chapter -- the same discipline
\texttt{docs/thesis/math\_manufacturing/thesis.tex} already applies to its
own Chapter~5.

\section{What changes when two source theses are read as one}


The two source theses divide the same subject matter along an ownership
line stated once, here, rather than re-argued in every chapter that could
plausibly claim a topic. \emph{The Projection Principle} owns the keystone
proper: the admission monoid, the manufacturing morphism $\muop$, the
receipt chain, and the Comprehension--Verification cost-regime gap
(Chapters~\ref{ch:keystone}--\ref{ch:gap} below), plus one chapter unique
to it, the trillion-agent instance (Chapter~\ref{ch:instance}).
\emph{Mission Geometry} owns the concrete planning and conformance
mechanics: PDDL grounding, the Petri-net marking polytope, Farkas
certificates, token-replay fitness, POWL separability, and cost
arbitration (Chapters~\ref{ch:pddl}--\ref{ch:cost} below). Where both
source outlines flag the same concept -- attention-as-capacity, the
Rice-quarantine rhetorical pattern, the cost vector -- this paper states
the general form once, in the chapter that owns it, and the concrete form
once, in the chapter that owns that, with an explicit cross-reference
rather than a restatement.


\begin{table}[h]
\centering
\small
\begin{tabular}{@{}p{4.2cm}p{4.2cm}p{4.2cm}@{}}
\toprule
\textbf{Shared concept} & \textbf{General form owned by} & \textbf{Concrete instantiation owned by} \\
\midrule
Attention as bounded capacity & Ch.~\ref{ch:attention}, \S\ref{sec:attn-general} (projection\_thesis) & Ch.~\ref{ch:attention}, \S\ref{sec:attn-pddl} (planning\_geometry) \\
Rice-quarantine retraction pattern & Ch.~\ref{ch:keystone}, \S\ref{sec:rice} (observation meanings) & Ch.~\ref{ch:separability} (workflow-net separability) \\
Cost vector / lexicographic order & Ch.~\ref{ch:cost}, \S\ref{sec:lex} (stated once, both sources verify it) & Ch.~\ref{ch:cost} (6-tuple \code{CostVector} operational form) \\
Marking polytope / Farkas / fitness / separability & Ch.~\ref{ch:petri}--\ref{ch:separability} (\emph{Mission Geometry} owns exposition) & \emph{projection\_thesis} cited only for the Rice-quarantine analogy tying separability to admission \\
\bottomrule
\end{tabular}
\caption{Ownership assignment for concepts flagged as shared by both source outlines. Stated once here; not re-argued per chapter.}
\end{table}


Two corrections are carried forward from the source material rather than
smoothed over. First, \emph{projection\_thesis}'s own formulation of the
manufacturing morphism is adopted exclusively: $\muop$ is deterministic
(M1) and bounded (M2) but need \emph{not} be injective, so Conservation of
Consequence (\S\ref{ch:keystone}) is stated as \emph{commitment}, not
semantic recovery, when $\muop$ is not injective -- fixing an earlier
naive-uniqueness reading the source outline flags explicitly. Second,
three of the keystone's central theorems -- Faithful Projection
(\texttt{thm:faithful}), the noncomprehension corollary
(\texttt{cor:noncomp}), and Conservation of Consequence
(\texttt{thm:conservation}) -- are, as of the receipts on disk, \emph{not}
machine-verified in either Lean lane: \texttt{thm:faithful} carries no
bare-core pilot receipt line at all and is recorded \texttt{excluded} in
the Mathlib-lane receipts; \texttt{cor:noncomp} and \texttt{thm:conservation}
are recorded \texttt{blocked} with zero attempts in both lanes. This paper
states all three as stated-but-not-machine-verified every place they are
used, including inside Chapter~\ref{ch:instance}'s instance argument,
which depends on \texttt{thm:faithful}.

\section{Structure of this document}


Chapter~\ref{ch:keystone} states the admission monoid, the Chatman
equation, and the receipt chain. Chapter~\ref{ch:gap} states the
verification-regime cost separation. Chapters~\ref{ch:pddl}--\ref{ch:cost}
give the planning-and-conformance geometry: PDDL grounding and bounded
branching, the attention balance law, the marking polytope, Farkas
separation certificates, token-replay fitness, POWL separability, and
lexicographic cost arbitration with the makespan functional.
Chapter~\ref{ch:instance} reports the trillion-agent status-byte instance.
Chapter~\ref{ch:implementation} merges the two sources' implementation-map
tables into one. An appendix (Chapter~\ref{ch:datalog}) reports the
Datalog index layer, demoted from a full chapter since it does not recur
elsewhere in either source. Chapter~\ref{ch:speculative} quarantines every
claim either source itself flags as an estimate, a consequence-not-theorem,
or a business-scale extrapolation, with one stated exception:
\texttt{thm:mrr} is placed there by this paper's own editorial judgment,
not by any hedge \emph{Mission Geometry} itself applies -- the source
states it as an unhedged theorem, verified in the bare-core Lean lane,
alongside \texttt{thm:lex}/\texttt{cor:dominate}/\texttt{def:makespan}
in Chapter~\ref{ch:cost}'s own source material
(\S\ref{ch:speculative} states the reclassification argument and its
receipt history in full). Chapter~\ref{ch:limitations} states
what is not shown, without euphemism, and Chapter~\ref{ch:conclusion}
closes.

\chapter{The Projection Keystone: Admission, Manufacture, and Receipts}
\label{ch:keystone}

\section{Rice quarantine and the admission map}
\label{sec:rice}


\begin{axiom}[Observation space]\label{ax:obs}
There is a set $\Obs$, the observation space, whose elements are
arbitrary finite records: logs, agent outputs, sensor frames, claims, tool
responses. $\Obs$ carries no decidable semantics -- ``does this
observation mean what it purports to mean'' is not assumed computable.
\end{axiom}
\begin{theorem}[Rice quarantine]\label{thm:rice}
Let $P$ be any non-trivial semantic property of the meanings observations
may encode. Then $\{o \in \Obs : P(o)\}$ is undecidable.
\end{theorem}
(\texttt{docs/thesis/rdf/projection\_thesis.ttl}, labels \texttt{ax:obs}/
\texttt{thm:rice}.) Both are verified in the bare-core pilot lane and,
since the migration to Mathlib, as a direct corollary of Mathlib's own
\texttt{ComputablePred.rice\_2} and \texttt{ComputablePred.halting\_problem}
(\texttt{Mathlib.Computability.Halting}, formalizing Carneiro 2019) rather
than a from-scratch hand proof -- a citation of an existing published
result, replacing 13 corpus-specific axioms in the bare-core file with
zero.


\begin{definition}[Admission map]\label{def:adm}
Fix a decidable $\Adm \subseteq \Obs$ on which a finite battery of
obligations $g_1,\dots,g_m : \Obs \to \{0,1\}$ is computable. The
admission map is the partial retraction $\adm : \Obs \rightharpoonup
\Adm \cup \{\Rfsl\}$, $\adm(o) = \rho(o) \in \Adm$ if $\bigwedge_i g_i(o)
= 1$, else $\Rfsl$, with $\rho$ a computable bounding/normalization and
$\Rfsl$ the distinguished refusal. Admission is idempotent on its image:
$\adm \circ \adm = \adm$.
\end{definition}
Rice quarantine (Theorem~\ref{thm:rice}) is the reason admission cannot
decide meaning in general; \texttt{def:adm} instead retracts onto a
decidable sub-language, verified in both Lean lanes
(\texttt{docs/thesis/rdf/projection\_thesis.ttl}, \texttt{def:adm}).

\section{The denial monoid}


\begin{construction}[Denial monoid]\label{con:denial}
Let $D = (\{0,1\}^n, \lor, \bm 0)$ be the commutative idempotent monoid of
denial words over $n$ independent obligation lanes, componentwise
\textsc{or}, identity $\bm 0$. Each obligation $g_i$ contributes a lane
map $d_i : \Obs \to \{0,1\}^n$ with $d_i(o) = \bm 0 \iff g_i(o) = 1$. The
total denial is $d(o) = \bigvee_i d_i(o)$, and $\adm(o) \ne \Rfsl \iff
d(o) = \bm 0$.
\end{construction}
\begin{proposition}[Semilattice, monotone admission]\label{prop:semilattice}
$(\{0,1\}^n, \lor)$ is the join-semilattice of $2^{[n]}$ with bottom
$\bm 0$. Adding obligations only enlarges denials: for $G' \supseteq G$,
$d_{G'}(o) \succeq d_G(o)$, so the admitted set shrinks monotonically
under $G' \supseteq G$.
\end{proposition}
Both verified in both Lean lanes (\texttt{con:denial}/\texttt{prop:semilattice},
\texttt{docs/thesis/rdf/projection\_thesis.ttl}). Refusal is data -- a bit
vector, not a side effect -- matching \texttt{CLAUDE.md}'s own invariant
that every error is a typed \texttt{Refusal} variant.


One further claim about the denial monoid appears only in
\emph{projection\_thesis}'s prose, in the source's own words: ``At a
trillion agents this taxonomy is a shared vocabulary: a refusal category
is legible to a recipient who never saw the sender's interior.'' The
source outline itself marks this claim \textsc{speculative} -- it is a
legibility/interpretive claim about a fixed enumeration of refusal
categories, not a corollary the denial-monoid construction above proves.
It is deferred to Chapter~\ref{ch:speculative} and is not asserted
here.

\section{The Chatman equation and the manufacturing morphism}


\begin{definition}[Manufacturing morphism, the Chatman equation]\label{def:mu}
$\muop : \Adm \to \Act$ is a deterministic computable map subject to (M1)
determinism/reproducibility -- $\muop(x)$ depends only on $x$, equal
admitted inputs give bit-identical artifacts -- and (M2) boundedness --
$\muop$ factors through a representation of size bounded by fixed
structural constants, so $\muop$ terminates with a priori bounded cost.
The Chatman equation is $\Act = \muop(\Adm)$, operationally $a =
\muop(\adm(o))$ when $\adm(o) \ne \Rfsl$, and $\muop(\Rfsl) = \Rfsl$.
\end{definition}
Verified in both Lean lanes (\texttt{def:mu},
\texttt{docs/thesis/rdf/projection\_thesis.ttl}). Determinism is not
injectivity: distinct admitted observations may manufacture an identical
artifact, which is exactly why Conservation of Consequence
(Theorem~\ref{thm:conservation}) must be stated in terms of commitment
rather than recovery.

\section{The law-bearing receipt}


\begin{definition}[Law-bearing receipt]\label{def:receipt}
Fix a collision-resistant hash $\chainH : \{0,1\}^{*} \to \{0,1\}^{256}$
and write $\dg(\cdot) = \chainH(\cdot)$. For a manufacture with admitted
observation $x = \adm(o)$, obligation set $G$ with denial word $d(o)$,
transition identity $\tau$, artifact $a = \muop(x)$, replay result
$\Fitness$, refusal reason $r$, and implementation version $v$, the frame
is $\mathsf{fr} = \langle \dg(x), \dg(G), d(o), \tau, \dg(a), \Fitness, r,
v \rangle$, and the chained receipt advances by $h_{+} = \chainH(h_{-}
\Vert \mathsf{fr})$. The receipt projection is $\Proj : \Sigma \to
\Recs$, $\Proj(\sigma) = (\text{verdict}, h_{+}, \Fitness, r)$, a tuple
of at most $\CL$ chunks.
\end{definition}
Verified in both Lean lanes (\texttt{def:receipt}). A receipt is a chained
commitment to named fields, never a bare byte-hash -- the same discipline
this project's own receipt tooling implements outside this thesis series
(\texttt{docs/thesis/02\_receipt\_cryptography.tex}).

\section{Faithful projection and its verification status}
\label{sec:faithful}


\begin{theorem}[Faithful projection]\label{thm:faithful}
The receipt is faithful to lawful consequence exactly to the extent the
frame commits its fields: for any committed field $\phi \in \{\dg(x),
\dg(G), d(o), \tau, \dg(a), \Fitness, r, v\}$, producing a trace
$\sigma'$ that differs in $\phi$ yet matches the claimed $h_+$ requires a
collision of $\chainH$. Fields not committed by the frame are not
attested.
\end{theorem}
This paper states Theorem~\ref{thm:faithful} but does not claim it is
machine-verified: \texttt{tools/paper-factory/lean-pilot/formalization\_receipts.jsonl}
carries no receipt line for \texttt{thm:faithful} at all (it is excluded a
priori from the bare-core run, since bare Lean~4 core has no library
support for its cryptographic/probabilistic content), and
\texttt{tools/paper-factory/lean-lake/mathlib\_migration\_receipts.jsonl}
records its status as \texttt{excluded} as well, with zero attempts. It is
reported here as \emph{stated, not machine-verified} in either lane, not
downgraded further nor upgraded to a proved status it does not have on
disk.


\begin{corollary}[Non-comprehension of boundary trust]\label{cor:noncomp}
An accepting verifier's predicate factors as $\mathrm{Verify} = V \circ
\Proj$ for a fixed $V : \Recs \to \{0,1\}$ of input size $\le \CL$.
Trust in the accept is a function of the $\CL$-sized projection and the
soundness of $V$, not of any agent's comprehension of $\sigma$.
\end{corollary}
Recorded \texttt{blocked} with zero attempts in both Lean lanes
(\texttt{cor:noncomp}). Stated, not machine-verified.

\section{Conservation of consequence}


\begin{theorem}[Conservation of Consequence]\label{thm:conservation}
(i) \emph{Causation}: every actuated artifact $a$ is the image of at
least one admitted observation, $a \in \im\muop$, and no $a \notin
\im\muop$ is actuated. (ii) \emph{Position uniqueness}: each actuation
appends exactly one receipt-chain position; the map $\text{actuation}
\mapsto h_+$ is injective up to hash collision. (iii) \emph{Observation
recoverability / commitment}: if $\muop$ is injective on $\Adm$, the
artifact determines the admitted observation (recoverability); otherwise,
when the frame commits $\dg(x)$, the receipt position uniquely commits to
$(\dg(x), \dg(G), \tau, \dg(a))$ (commitment, not semantic recovery).
\end{theorem}
Recorded \texttt{blocked} with zero attempts in both Lean lanes
(\texttt{thm:conservation}). This is the corrected formulation
(\S\ref{sec:thesis-statement} above): clause (iii) is deliberately weaker
than a naive uniqueness claim would be, precisely because $\muop$ is only
known deterministic, not injective. Being unattempted in Lean is a scope
gap, not a proof failure; it is reported as such, not conflated with a
rejected kernel run.

\chapter{Cost Regimes and the Comprehension--Verification Gap}
\label{ch:gap}

\section{Verification regimes}


\begin{definition}[Verification regimes]\label{def:regimes}
For a trace of $n$ committed frames: boundary projection inspection (any
verifier, per message) costs $O(\CL)$; single-frame recomputation (spot
check) costs $O(1)$ per frame; checkpointed/indexed replay (sampled audit)
costs $O(\log n)$ or a bounded path; full chain replay (full audit) costs
$O(n)$.
\end{definition}
Verified in both Lean lanes (\texttt{def:regimes},
\texttt{docs/thesis/rdf/projection\_thesis.ttl}). This table is stated
once, here; Chapters~\ref{ch:conformance} and \ref{ch:fitness} cite it
rather than restate it when discussing certificate- and
replay-verification cost.

\section{The Fleet Trust Budget}
\label{sec:fleet-budget}


\begin{theorem}[Comprehension--Verification Gap; Fleet Trust Budget]\label{thm:gap}
Let $\mathrm{Comp}(\sigma) = \Theta(\dim\Sigma)$ be the cost of
comprehending the interior and $\mathrm{Ver}_{\partial}(\sigma) =
O(\CL)$ the cost of boundary verification of $\Proj(\sigma)$. Then
$\mathrm{Comp}(\sigma)/\mathrm{Ver}_{\partial}(\sigma) =
\Omega(\dim\Sigma/\CL) \to \infty$ as $\dim\Sigma \to \infty$, while
boundary verification stays constant. For a fleet, boundary cost is
$O(M \cdot \CL)$ for $M$ messages, $O(N \cdot \CL)$ per round,
$O(N^2 \cdot \CL)$ all-to-all -- linear in $\CL$, not interior dimension.
\end{theorem}
Verified in both Lean lanes (\texttt{thm:gap}). This is \emph{not} a claim
that all mechanical audit is $O(1)$: full replay is honestly $O(n)$
(Definition~\ref{def:regimes}); the constant-cost claim is confined to
boundary inspection of an already-committed projection, a distinction
\emph{projection\_thesis} itself insists on and this paper preserves.

\chapter{PDDL Action Law: Lifted Domains, Grounding, Bounded Branching}
\label{ch:pddl}

\section{Lifted domains and grounding}


\begin{definition}[Lifted domain]\label{def:domain}
A lifted domain is a tuple $\mathcal{D} = (\mathcal{T}, \mathcal{P},
\mathcal{F}, \mathcal{S})$: a finite type hierarchy $\mathcal{T}$, a
finite set $\mathcal{P}$ of predicate symbols, a finite set $\mathcal{F}$
of numeric function symbols (fluents), and a finite set $\mathcal{S}$ of
action schemas; a durative schema carries typed parameters, a duration
constraint, a condition $C_s$, and an effect $E_s$.
\end{definition}
Verified in both Lean lanes (\texttt{def:domain},
\texttt{docs/thesis/rdf/03\_planning\_geometry.ttl}).


\begin{definition}[Grounding]\label{def:ground}
Given a domain $\mathcal{D}$ and a problem $(\mathcal{Ob}, \bm m_0,
\nu_0, \gamma)$, grounding substitutes every type-compatible object
tuple into each schema's parameters, producing a finite set of ground
actions; a grounded state is a pair $(\bm m, \nu)$ of a set $\bm m$ of
true ground atoms and a fluent valuation $\nu$.
\end{definition}
Verified in both Lean lanes (\texttt{def:ground}).

\section{Bounded branching}
\label{sec:bounded-ground}


\begin{theorem}[Bounded branching]\label{thm:bounded-ground}
Let $\mathcal{D}$ have schemas of maximum parameter arity $k$ over
$|\mathcal{Ob}| = N$ objects. The number of ground actions is at most
$|\mathcal{S}| \cdot N^{k}$, a finite constant independent of the goal;
consequently a forward search over grounded actions has a branching
factor bounded a priori, and $\muop$'s manufacture terminates within a
fixed depth bound.
\end{theorem}
Verified in the bare-core Lean lane
(\texttt{tools/paper-factory/lean-pilot/formalization\_receipts.jsonl},
\texttt{thm:bounded-ground}, 1 attempt). In the Mathlib-linked lane it is
recorded \texttt{unformalized} after 3 attempts, a regression against the
bare-core proof: the final attempt fails on \texttt{Unknown constant
\`{}Finset.sum\_le\_sum\`{}} and \texttt{Unknown identifier
\`{}smul\_eq\_mul\`{}} against the current Mathlib snapshot
(\texttt{tools/paper-factory/lean-lake/mathlib\_migration\_receipts.jsonl})
-- a lookup failure against a moving library target, not a defect in the
bare-core proof, which remains verified on its own terms.


In \code{bcinr-pddl} -- a sibling git repository at \texttt{../bcinr},
outside praxis's own \texttt{crates/}/\texttt{src/} trees, reached only via
the path dependency in \texttt{Cargo.toml}, not vendored into this
codebase -- the ground-action count is additionally hard-capped
by \code{PDDL8\_MAX\_GROUND}; exceeding it is a typed \code{BoundExceeded}
refusal rather than an exception or unbounded search, and its
\code{GroundProblem::find\_plan} (\texttt{bcinr-pddl/src/ground.rs}) is
breadth-first search over ground-atom sets whose own failure mode is a
receipted refusal (\code{NoAdmittedPlan}), never a divergence. Praxis's own
in-repo \texttt{crates/pddl-index/src/ground.rs} implements the analogous
but distinctly named \code{IndexedGroundProblem::find\_plan}, over its own
\code{GroundError} enum; the two are separate types in separate crates and
this paragraph cites only the former, \code{bcinr-pddl}'s. This is a
code-level remark, not a separate corpus statement, cited in the source
outline against no additional Lean label.

\section{Relational membership via XOR filters}


\begin{construction}[Relational membership via XOR filter]\label{con:xorf}
Let $R \subseteq \Obs$ be the reachability fixpoint fact set. A frozen
fact store builds an 8-bit XOR filter over FNV-1a hashes of $R$, with
index mapping via the fastrange reduction $\mathrm{reduce}(x,n) = (x
\times n)/2^{32}$; membership is gated by $x \in B$ if $F(x) =
\mathrm{true}$, else false, with false-positive rate $\approx 0.4\%$ and
false negatives $0$.
\end{construction}
Verified in both Lean lanes (\texttt{con:xorf}).

\chapter{Numeric Fluents and the Attention Balance Law}
\label{ch:attention}

\section{The general attention calculus}
\label{sec:attn-general}


\begin{definition}[Attention balance, general form]\label{def:attnbalance}
Let $c(t) \ge 0$ be capacity and admitted action $j$ draw rate $r_j \ge
0$ on $[s_j, e_j]$. Free capacity is $f(t) = c(t) - \sum_{j : t \in
[s_j,e_j]} r_j$. A schedule is feasible iff $f(t) \ge 0$ for all $t$.
\end{definition}
Verified in the bare-core Lean lane
(\texttt{docs/thesis/rdf/projection\_thesis.ttl}, \texttt{def:attnbalance});
recorded \texttt{unformalized} in the Mathlib lane after a Mathlib
big-operator syntax mismatch across 3 attempts (budget exhausted, per the
source outline). This is the general statement of attention-as-bounded-
capacity that \S\ref{sec:attn-pddl} below instantiates concretely for a
single PDDL fluent.


\begin{proposition}[Attention conservation, general form]\label{prop:attnconservation}
Total attention consumed, $\int_0^T \sum_j r_j \mathbf{1}_{[s_j,e_j]}(t)\,
dt = \sum_j r_j(e_j - s_j)$, is independent of schedule ordering.
Reordering redistributes \emph{when} attention is spent, not \emph{how
much}; the optimization target is the makespan functional $T[\cdot]$
(\S\ref{sec:makespan} below) under precedence and the pointwise capacity
constraint.
\end{proposition}
Verified in the bare-core Lean lane; recorded \texttt{blocked} in the
Mathlib lane (\texttt{prop:attnconservation}).

\section{The concrete PDDL fluent mechanics}
\label{sec:attn-pddl}


\begin{definition}[Fluent draw, PDDL instantiation]\label{def:balance}
A durative action $j$ draws on a fluent $\phi$ if its effect decreases
$\phi$ by $r_j$ at start and increases it by $r_j$ at end for rate $r_j
\ge 0$, holding $r_j$ units over $[s_j, e_j)$; the free level of $\phi$
at time $t$ is $f_\phi(t) = \nu_0(\phi) - \sum_{j : t \in [s_j,e_j)}
r_j$.
\end{definition}
Verified in both Lean lanes (\texttt{def:balance},
\texttt{docs/thesis/rdf/03\_planning\_geometry.ttl}). This is
Definition~\ref{def:attnbalance} specialized to a single PDDL numeric
fluent under the borrow/return action pattern.


\begin{proposition}[Conservation of the borrowed fluent]\label{prop:conserve}
If every action drawing on $\phi$ both decrements it by $r_j$ at start and
increments it by $r_j$ at end, the net effect of each completed action on
$\phi$ is zero, and the terminal valuation equals the initial:
$\nu_{\mathrm{end}}(\phi) = \nu_0(\phi)$; scheduling redistributes when
$\phi$ is held, never how much exists.
\end{proposition}
Verified in both Lean lanes (\texttt{prop:conserve}) -- attention is
redistributed in time, never created or destroyed.


\begin{proposition}[Feasibility as free-level nonnegativity]\label{prop:balance}
A durative-action schedule is feasible with respect to $\phi$ iff
$f_\phi(t) \ge 0$ for all $t$; for the unit-rate attention fluent this
says the number of concurrently in-flight capabilities never exceeds the
capacity $\nu_0(\code{attention})$.
\end{proposition}
Verified in both Lean lanes (\texttt{prop:balance}) -- exactly
Definition~\ref{def:attnbalance}'s general feasibility condition,
restated as the PDDL fluent's own guard.

\chapter{The Marking Polytope: Petri Nets, State Equation, Cone Reachability}
\label{ch:petri}

\section{Nets, incidence, and the state equation}


\begin{definition}[Net]\label{def:net}
A net has $p$ places and a finite set $T$ of transitions; a marking is a
vector $\bm m \in \Zplus^{p}$; a transition $t$ is a pair
$(\bm m^{-}_{t}, \bm m^{+}_{t})$ of preset and postset, enabled at
$\bm m$ iff $\bm m \ge \bm m^{-}_{t}$ coordinatewise, firing to
$\bm m' = \bm m - \bm m^{-}_{t} + \bm m^{+}_{t}$.
\end{definition}
Verified in both Lean lanes (\texttt{def:net},
\texttt{docs/thesis/rdf/03\_planning\_geometry.ttl}).


\begin{definition}[Incidence matrix, state equation]\label{def:incidence}
Let $\Nmat \in \mathbb{Z}^{p \times |T|}$ have column $\bm\delta_t =
\bm m^{+}_{t} - \bm m^{-}_{t}$ for each transition $t$; a firing sequence
with Parikh vector $\bm x \in \Zplus^{|T|}$ satisfies the state equation
$\bm m = \bm m_0 + \Nmat\bm x$.
\end{definition}
Verified in both Lean lanes (\texttt{def:incidence}).

\section{The reachability cone and its polytope relaxation}


\begin{proposition}[Cone necessity]\label{prop:cone}
Every marking $\bm m$ reachable from $\bm m_0$ satisfies the state
equation for some $\bm x \in \Zplus^{|T|}$, hence lies in the integer
points of $\bm m_0 + \cone(\Nmat)$ intersected with $\Zplus^{p}$; the
state equation is necessary for reachability but not sufficient.
\end{proposition}
Recorded \texttt{unformalized} in the bare-core Lean lane after 3
attempts: the induction step fails on an index-shift mismatch in
\code{incidenceSumAux} and an \code{omega} failure relating a \code{Nat}
truncated-subtraction expression to its \code{Int}-cast components
(\texttt{tools/paper-factory/lean-pilot/formalization\_receipts.jsonl},
\texttt{prop:cone}). Verified in the Mathlib-linked lane in 2 attempts
(\texttt{tools/paper-factory/lean-lake/mathlib\_migration\_receipts.jsonl}),
depending only on Lean's three standard foundational axioms. The
insufficiency half (spurious solutions exist) is the classical Petri-net
fact of Murata 1989, cited rather than reproved in either lane -- an
external-citation proof step, not an in-paper derivation.


\begin{definition}[Marking polytope]\label{def:polytope}
The marking polytope of a net is the relaxation $\Poly = \{\bm m_0 +
\Nmat\bm x : \bm x \ge \bm 0\} \cap \{\bm m \ge \bm 0\} \subseteq
\Rplus^{p}$, the continuous translated cone whose integer points
over-approximate the reachable set.
\end{definition}
Recorded \texttt{blocked} (not attempted) in the bare-core lane, verified
in the Mathlib-linked lane (\texttt{def:polytope}). \emph{projection\_thesis}'s
own \texttt{def:reachcone} states the same reachable-set-in-a-cone content
in different notation; this paper states it once, here, as
\texttt{def:polytope}, and does not restate \texttt{def:reachcone}
separately.

\section{STRIPS grounding as a place/transition net}


\begin{construction}[STRIPS as a place/transition net]\label{con:strips}
Take the places to be the ground atoms; a ground action $t$ with
delete-effects $\bm m^{-}_{t}$ and add-effects $\bm m^{+}_{t}$ is a
transition, and \code{GroundProblem::find\_plan} searches marking space
by BFS with the successor relation exactly transition firing.
\end{construction}
\begin{proposition}[Safe-net bit identity]\label{prop:safe}
On a safe (1-bounded) net, the integer firing rule $\bm m' = \bm m -
\bm m^{-}_{t} + \bm m^{+}_{t}$ coincides bit-for-bit with the branchless
bitset update \code{enabled\_tokens} $\gets$ (\code{enabled\_tokens} \&
$\lnot\bm m^{-}_{t}$) $|$ $\bm m^{+}_{t}$, and the enabling test coincides
with the branchless subset check (\code{enabled} \& $\bm m^{-}_{t}$)
$\oplus \bm m^{-}_{t} = 0$.
\end{proposition}
Both verified in both Lean lanes (\texttt{con:strips}, \texttt{prop:safe}).
\code{PowlReplayVerifier}'s branchless update is thus not an approximation
of the Petri-net firing rule on safe nets -- it is the same integer
operation, proved coincident.

\chapter{Conformance as Cone Membership: Farkas Separation Certificates}
\label{ch:conformance}

\section{Conformance as polytope membership}


\begin{definition}[Conformance]\label{def:conf}
A trace with final marking $\bm m^{\dagger}$ and transition-count vector
$\bm x$ conforms to a net with initial marking $\bm m_0$ if $\bm
m^{\dagger} \in \Poly$ and the $\bm x$ witnessing it is realized by a
genuine firing ordering; it fails membership if $\bm m^{\dagger} \notin
\Poly$.
\end{definition}
Recorded \texttt{blocked} (not attempted) in the bare-core Lean lane;
verified in the Mathlib lane (\texttt{def:conf}). \emph{projection\_thesis}'s
\texttt{prop:conformembership} states the same polytope-membership
content; it is not restated separately here.

\section{Farkas' lemma and the separating certificate}


\begin{theorem}[Farkas separation]\label{thm:farkas}
Fix $\bm b = \bm m^{\dagger} - \bm m_0$. Exactly one holds: (a) there
exists $\bm x \ge \bm 0$ with $\Nmat\bm x = \bm b$; or (b) there
exists $\bm y \in \mathbb{R}^{p}$ with $\Nmat^{\top}\bm y \le \bm 0$
and $\bm y^{\top}\bm b > 0$, a separating hyperplane certifying
$\bm m^{\dagger} \notin \bm m_0 + \cone(\Nmat)$, and if (b) holds no
firing sequence reaches $\bm m^{\dagger}$.
\end{theorem}
Recorded \texttt{blocked} (not attempted) in the bare-core Lean lane.
Recorded \texttt{unformalized} in the Mathlib lane after 1 attempt, with
the exact missing lemma named: \code{PointedCone.FG.isClosed} does not
exist in the current Mathlib snapshot. Mathlib's topological Farkas
machinery (\code{ProperCone.hyperplane\_separation\_point},
\texttt{Mathlib.Analysis.Convex.Cone.InnerDual}) would give the
alternative once $\cone(\Nmat)$ is packaged as a \code{ProperCone}, but
that packaging requires \code{IsClosed}$(\cone(\Nmat))$ -- the
topological half of the Minkowski--Weyl theorem (a finitely generated cone
in a finite-dimensional real vector space is closed) -- which the current
Mathlib snapshot does not supply as a reusable lemma. This is reported
\texttt{unformalized}, not faked with \code{sorry} or an axiom. The
proof-in-paper itself leans on Farkas' lemma as an external citation
(Schrijver 1986), not an in-paper derivation, for its own first step.


\emph{projection\_thesis}'s own \texttt{prop:conformembership} states
materially the same claim -- conformance is polytope membership,
nonconformance yields a Farkas certificate -- and is recorded
\texttt{unformalized} in the Mathlib lane for a \emph{different} reason: a
\code{Matrix.dotProduct} identifier mismatch against Mathlib's unqualified
\code{dotProduct}, budget exhausted after 3 attempts. Both labels name the
same mathematical content under two independent Lean files with two
independent kernel-error causes; this paper cites \texttt{thm:farkas} as
the canonical statement (owned by \emph{Mission Geometry}) and records
\texttt{prop:conformembership}'s distinct failure mode here rather than
silently merging the two into one reported status.

\section{The certificate's verification cost}


\begin{corollary}[Nonconformance certificate]\label{cor:cert}
A nonconformance certificate is a single vector $\bm y \in
\mathbb{R}^{p}$ together with $\Nmat^{\top}\bm y \le \bm 0$ and
$\bm y^{\top}\bm b > 0$; checking it is a matrix product and two sign
tests, independent of trace length, verifiable within a bounded context
$\CL$.
\end{corollary}
Recorded \texttt{blocked} (not attempted) in both Lean lanes
(\texttt{cor:cert}). Its cost bound, $O(p \cdot |T|)$, is exactly the
boundary-verification regime of Definition~\ref{def:regimes}: a
nonconformance certificate is checkable at the same $O(\CL)$-shaped cost
as any other receipted claim, independent of the trace's own length.

\chapter{Token-Replay Fitness as Normalized Distance}
\label{ch:fitness}

\section{The replay verifier}


\begin{definition}[Replay verifier]\label{def:replay}
These field names are quoted verbatim from the \code{PowlReplayVerifier}
struct (\texttt{bcinr-powl-receipt/src/replay.rs:31--35}), an external
crate reached only via the path dependency declared in
\texttt{Cargo.toml} -- a sibling git repository at \texttt{../bcinr},
outside praxis's own \texttt{crates/}/\texttt{src/} trees; praxis's own
\texttt{crates/praxis-core/src/replay\_adapter.rs} only imports and wraps
\code{bcinr\_powl\_receipt::replay::PowlReplayVerifier} and does not
itself define the type. The verifier holds a marking $\bm m$
(\code{enabled\_tokens}) initialized
to the entry token, and accumulators \code{replayed}, \code{fitted},
\code{enabled\_not\_taken}; replaying a frame rejects an invalid node bit,
requires $\bm m \ge \bm m^{-}$, records enabled-but-unconsumed tokens,
fires $\bm m \gets (\bm m \,\&\, \lnot\bm m^{-}) \,|\, \bm m^{+}$, and
sets the node bit in \code{replayed} and \code{fitted}.
\end{definition}
Verified in both Lean lanes (\texttt{def:replay},
\texttt{docs/thesis/rdf/03\_planning\_geometry.ttl}).


\begin{theorem}[Replay dichotomy]\label{thm:fitness}
Replaying a sequence of frames against the net either (i) completes with
no violation, in which case it is a genuine firing sequence and $\Fitness
= 1$; or (ii) halts at the first frame whose preset is not contained in
the current marking, returning a coordinate-localized witness of the
earliest disabled transition.
\end{theorem}
Verified in both Lean lanes (\texttt{thm:fitness}).

\section{Fitness and precision as population ratios}


\begin{definition}[Fitness and precision]\label{def:fitness}
On finalization the verifier returns, in Q16.16 fixed point, $\Fitness =
|\code{fitted}|/|\code{replayed}|$ and precision $=
|\code{replayed}|/(|\code{replayed}| + |\code{enabled\_not\_taken}|)$,
where $|\cdot|$ is population count and $\Fitness \coloneqq 1$ when the
denominator is zero.
\end{definition}
\begin{proposition}[Honest, orthogonal ratios]\label{prop:honest}
$\Fitness$ and precision are ratios of bit-populations of
disjointly-maintained bitsets; neither can exceed $1$, and $\Fitness = 1$
is attained exactly on completed replays; \code{enabled\_not\_taken} feeds
precision, not fitness, so the two metrics measure orthogonal deviations
and cannot be traded against each other.
\end{proposition}
Both verified in both Lean lanes (\texttt{def:fitness}, \texttt{prop:honest}).
\emph{projection\_thesis}'s \texttt{def:fitnessdistance}/\texttt{thm:unitfitness}
restate the same fitness-as-normalized-distance content ($\Fitness = 1 -$
the fraction of tokens forced onto unenabled transitions) and are verified
in the bare-core lane, blocked in the Mathlib lane; this paper treats
\texttt{def:fitness}/\texttt{thm:fitness} above as the canonical statement
and does not duplicate the definition a second time.

\chapter{POWL 2.0: Recursive Partial Orders and Choice Graphs}
\label{ch:powl}

\section{The POWL grammar}


\begin{definition}[POWL model]\label{def:powl}
A POWL model over an activity alphabet $A$ is one of: an activity $a \in
A$ (a leaf); a partial order $(\{M_1,\dots,M_k\}, \prec)$ executed by
respecting $\prec$ and running incomparable children concurrently; or a
choice graph $(\{M_1,\dots,M_k\}, E)$ combined by exclusive choice (and
cyclic/loop edges, POWL 2.0), taking exactly one live branch.
\end{definition}
Verified in both Lean lanes (\texttt{def:powl},
\texttt{docs/thesis/rdf/03\_planning\_geometry.ttl}).


\begin{proposition}[Bounded local geometry]\label{prop:local}
In a POWL model, the token marking local to a node ranges over the tokens
of that node's immediate children: a node of arity $k$ has a local marking
space of dimension $O(k)$, independent of the total size of the model; a
node's replay is checked in a working set bounded by its arity.
\end{proposition}
Verified in both Lean lanes (\texttt{prop:local}).

\section{Temporal plans as Hasse diagrams}


\begin{construction}[Temporal plan to POWL tape]\label{con:tape}
\code{temporal\_plan\_to\_powl\_tape} (\texttt{bcinr-pddl/src/powl\_bridge.rs},
a sibling git repository at \texttt{../bcinr}, outside praxis's own
\texttt{crates/}/\texttt{src/} trees, reached only via the path
dependency declared in \texttt{Cargo.toml} -- not the in-repo
\texttt{crates/powl2-decompose} crate, which has no tape/mask machinery)
converts a \code{TemporalPlan} into a
tape of \code{PowlOpSpec}s (struct also defined in that same external
file): each step $i$ becomes an activity with a
\code{pred\_mask} of steps $j<i$ finishing before $i$ starts and a one-hot
\code{succ\_mask} (both fields on the tape type in
\texttt{bcinr-powl/src/tape.rs}, same external repository); a second pass
performs transitive reduction, clearing
bits of predecessors' own predecessors, leaving only direct predecessors.
\end{construction}
\begin{proposition}[Hasse diagram]\label{prop:hasse}
Let $\prec$ be the strict order ``$j$ finishes before $i$ starts'' on
plan steps. The initialized \code{pred\_mask} of step $i$ is $\{j : j
\prec i\}$, and after transitive reduction it is the covering relation of
$\prec$; the resulting dependency graph is a DAG, and two steps are
schedulable concurrently iff $\prec$-incomparable.
\end{proposition}
Both verified in both Lean lanes (\texttt{con:tape}, \texttt{prop:hasse}).

\chapter{Separability as a Process-Layer Rice Quarantine}
\label{ch:separability}

\section{The decomposition theorem}


\begin{theorem}[Separable decomposition]\label{thm:sep}
A sound, separable workflow net admits a recursive decomposition into a
POWL model whose every internal node is a partial order or a choice
graph and whose leaves are activities, with each node's local marking
geometry dimension bounded by its arity, independent of the global net
size.
\end{theorem}
Verified in both Lean lanes (\texttt{thm:sep},
\texttt{docs/thesis/rdf/03\_planning\_geometry.ttl}). The underlying
decomposition result is attributed to Kourani and van der Aalst and
cited, not reproved, in the source paper itself; the Lean formalization on
record should be read as capturing the structural corollary this paper
states -- bounded local marking dimension under decomposition -- not as an
independent re-derivation of the cited external decomposition theorem.
\emph{projection\_thesis} restates the same theorem as its own
\texttt{thm:sep} (verified both lanes) and additionally frames
separability as ``a process-layer Rice quarantine''; that framing is the
general Rice-quarantine pattern of \S\ref{sec:rice}, specialized to
processes, and is stated here as the specialization, not re-derived.

\section{Separability as a decidable admission predicate}


\begin{proposition}[Separability as a decidable admission predicate]\label{prop:quarantine}
Let $\mathrm{sep}$ be the decidable predicate ``the net is sound and
separable.'' It is a legitimate admission obligation: total and
computable, retracting the space of arbitrary control-flow onto the
decidable sub-language of POWL-expressible processes; a separable net is
admitted, a non-separable net is refused with reason.
\end{proposition}
Verified in both Lean lanes (\texttt{prop:quarantine}). This is the
process-layer instance of Definition~\ref{def:adm}'s admission map, with
$\mathrm{sep}$ as the obligation and the POWL decomposition as $\rho$.

\section{Normalization and language preservation}


\begin{construction}[Sub-net normalization]\label{con:normalize}
Given a partition part $T' \subseteq T$ with entry places
$P_{\mathrm{in}}$ and exit places $P_{\mathrm{out}}$, the sub-net is
normalized to a valid WF-net by adding fresh boundary places $p_s, p_e$
and silent transitions $\tau_{\mathrm{in}}, \tau_{\mathrm{out}}$,
redirecting boundary arcs onto a unified source and sink.
\end{construction}
\begin{theorem}[Language-preserving conversion]\label{thm:lang-correct}
Let $N$ be a safe and sound workflow net, and let $\psi =
\mathrm{convert}(N)$ be its POWL~2.0 conversion. Then for any prefix
trace length $k$, $\mathcal{L}_k(N) = \mathcal{L}_k(\psi) =
\mathcal{L}_k(\mathrm{recompose}(\psi))$.
\end{theorem}
\texttt{con:normalize} is verified in both Lean lanes.
\texttt{thm:lang-correct} carries no proof body at all in the source
\texttt{.tex} -- no \code{begin\{proof\}} block -- and is reported here as
asserted, not derived, in the source paper, even though the corpus
records a Lean formalization for it in both lanes; what exactly that Lean
statement encodes should be checked against
\texttt{tools/paper-factory/lean-pilot/thm\_lang\_correct.lean} before
citing it as a fully in-paper-proved result, a distinction this paper
flags rather than smooths over.

\chapter{Cost Arbitration: Lexicographic Order and the Makespan Functional}
\label{ch:cost}

\section{The cost vector}


\begin{definition}[Cost vector]\label{def:cost}
Assign a plan $\bm c = (c_0, c_1, \dots, c_k)$, with $c_0 \in \{0,1\}$
the unadmitted indicator ($0=$ admitted) and $c_1,\dots,c_k$ bounded
secondary costs, ordered lexicographically ($\preceq_{\mathrm{lex}}$),
cheapest-first. The router's operational instance is the 6-tuple
\code{CostVector} $= (c_0,\dots,c_5) = (\code{admitted},
\code{unreceipted\_mutation\_risk}, \code{human\_attention\_seconds},
\code{token\_cost}, \code{latency\_ms}, \code{context\_switches})$ (field
names verbatim from the \code{CostVector} struct at
\texttt{capability\_router.rs:107--122} in the \code{bcinr-pddl} crate --
a sibling git repository at \texttt{../bcinr}, outside praxis's own
\texttt{crates/}/\texttt{src/} trees, reached only via the path
dependency declared in \texttt{Cargo.toml}; this and the other
\code{bcinr-*} code loci cited in this chapter and Chapter~\ref{ch:powl}
(\texttt{con:tape}) are not part of the reviewed praxis codebase itself),
with the refused vector \code{refused()} $=
(\text{unadmitted}, 255, \infty, \infty, \infty, 255)$ as the top
element.
\end{definition}
Both the general form (\texttt{def:costvector}, \emph{projection\_thesis})
and the 6-tuple operational form (\texttt{def:cost}, \emph{Mission
Geometry}) are verified in both Lean lanes. They are stated once here as
one definition, general form first, operational instance second, rather
than twice.

\section{Lexicographic order as an infinite-weight limit}
\label{sec:lex}


\begin{theorem}[Lexicographic order as an infinite-weight limit]\label{thm:lex}
Let secondary coordinates be bounded, $|c_i| \le M$ for $i \ge 1$. Define
$C_\lambda(\bm c) = \sum_{i=0}^{k} \lambda^{k-i} c_i$ with $\lambda >
1$. There is a threshold $\Lambda$ such that for all $\lambda >
\Lambda$, $\bm c \preceq_{\mathrm{lex}} \bm c' \iff C_\lambda(\bm c)
\le C_\lambda(\bm c')$; as $\lambda \to \infty$ the weight on $c_0$
diverges relative to all others.
\end{theorem}
\begin{corollary}[Unlawfulness dominates]\label{cor:dominate}
An admitted plan sorts strictly before any refused plan regardless of its
secondary costs: for any finite risk, makespan, token, latency, and switch
values, an admitted \code{CostVector} is $\prec_{\mathrm{lex}}$ the
refused top element; the price of unlawfulness is infinite in the limit
that defines the order.
\end{corollary}
Both verified in both Lean lanes, in both source labelings
(\texttt{thm:lex}/\texttt{cor:dominate}, \emph{Mission Geometry};
\texttt{thm:lex}, \emph{projection\_thesis}). Routing determinism is
receipted -- identical tasks produce identical BLAKE3 route chains -- a
code-level remark citing the
\code{routing\_is\_deterministic\_same\_task\_same\_route} test
(\texttt{bcinr-pddl/src/capability\_router.rs:340}, an admitted two-effect
task asserting equal \code{route\_chain} across two routing calls), with
no separate corpus label.

\section{The makespan functional}
\label{sec:makespan}


\begin{definition}[Makespan functional]\label{def:makespan}
Given the precedence DAG with per-op durations $d_i$ and release times,
the forward pass computes $\mathrm{es}_i = \max(r_i, \max_{j \prec i}
\mathrm{ef}_j)$, $\mathrm{ef}_i = \mathrm{es}_i + d_i$; the makespan is
$T = \max_i \mathrm{ef}_i$; the backward pass computes $\mathrm{ls}_i =
\mathrm{lf}_i - d_i$; the slack of op $i$ is $\mathrm{ls}_i -
\mathrm{es}_i$, and the critical path is the set of zero-slack ops.
\end{definition}
\begin{proposition}[Longest path]\label{prop:critical}
The makespan $T$ equals the length of the longest ($\prec$-)path in the
precedence DAG weighted by durations, and an op has zero slack iff it lies
on some longest path.
\end{proposition}
\begin{proposition}[Capacity sensitivity]\label{prop:sensitivity}
Let $T(\phi = c)$ be the makespan the planner finds when fluent $\phi$'s
initial capacity is $c$. The capacity delta reports $(T(c-1), T(c),
T(c+1))$, and $\phi$ is binding iff $T(c-1) \ne T(c)$ (or infeasible at
$c-1$); this is a finite-difference sensitivity of the specific schedule
the greedy planner produced, not a dual shadow price of an optimal
scheduler -- the source paper explicitly disclaims the stronger reading.
\end{proposition}
All three verified in both Lean lanes (\texttt{def:makespan},
\texttt{prop:critical}, \texttt{prop:sensitivity}).

\chapter{Instance: An Autonomous Build and the Trillion-Agent Status Byte}
\label{ch:instance}

\section{An unread interior, trusted by boundary receipts}


An autonomous build with an interior of roughly $2.5 \times 10^6$ tokens,
unread by any agent, is trusted entirely via boundary receipts: verdict,
chain hash, fitness $=1.0$, receipted refusal count, pass rate, and an
oracle-agreement flag. This claim, as \emph{projection\_thesis} states it,
cites Theorem~\ref{thm:faithful} and Theorem~\ref{thm:gap} together with
Proposition~\ref{prop:semilattice}; it therefore inherits
Theorem~\ref{thm:faithful}'s unresolved verification status
(Ch.~\ref{ch:keystone} above) exactly, and this paper states the
inheritance rather than silently upgrading the build-trust claim past what
\texttt{thm:faithful} itself has earned on disk.

\section{The eight-bit status byte}


\begin{construction}[Eight-bit status byte]\label{con:agent8}
Project each agent to a status byte $b \in \{0,1\}^8$ whose lanes name
admitted, evidence-ok, within-budget, authority-bound, healthy,
conformant, receipted, replayable. A fleet of $N$ agents is a bit-packed
vector; a 64-bit word holds 8 agents; a fleet admission sweep is the
branchless mask reduction of the denial construction (Construction~\ref{con:denial})
applied word-parallel.
\end{construction}
Verified in both Lean lanes (\texttt{con:agent8}). Each bit is set only by
a named mechanism, never by self-report -- the worked numeric example in
the source: a full status byte \code{0xFF} when every lane passes,
dropping to \code{0x3F} with refusal category ``Prerequisites'' when
ledger evidence is absent.

\section{Fleet arithmetic}


\begin{proposition}[Fleet arithmetic, exact count]\label{prop:planetary}
For a fleet of $N$ agents under Construction~\ref{con:agent8}: the status
vector occupies exactly $N$ bytes, packed to $N/8$ 64-bit words; a full
admission sweep is $N/8$ masked-\textsc{or} word operations.
\end{proposition}
This part of \texttt{prop:planetary} (part (a) in the source) is exact
counting, verified in both Lean lanes, and is the only part of
\texttt{prop:planetary} stated in this chapter. Parts (b)--(d) -- the
compressed-status estimate, the affordability-ratio estimate, and the
consequence-not-theorem about planetary control -- are engineering and
order-of-magnitude estimates the source itself labels as such, and are
quarantined into Chapter~\ref{ch:speculative} rather than stated here
alongside the exact count.

\section{Five named adversaries}


The source instance names five adversaries -- forged self-report, broken
replay, missing policy digest, non-separable workflow, unreceipted claim
-- each defeated by a specific mechanism rather than by inspecting the
interior: forged self-report by Theorem~\ref{thm:faithful} (stated, not
machine-verified, Ch.~\ref{ch:keystone}), broken replay by
Theorem~\ref{thm:fitness} (verified, Ch.~\ref{ch:fitness}; the
\emph{projection\_thesis} corpus additionally labels the same content
\texttt{thm:unitfitness}), non-separable workflow by Theorem~\ref{thm:sep}
(verified, Ch.~\ref{ch:separability}), and unreceipted claim by
\texttt{prop:selfauditedscale}, which is itself recorded \texttt{blocked}
with zero attempts in both Lean lanes and is quarantined to
Chapter~\ref{ch:speculative} below as the ``antibody clause.'' This
chapter's adversary list therefore inherits a mixed evidentiary status
across its five entries, stated per-entry rather than as one uniform
``defeated'' claim.

\chapter{Implementation Map}
\label{ch:implementation}


Both source theses carry their own implementation-map appendix,
structurally identical in form (a table of formal object $\to$ code
locus); consolidating them into two merged tables here replaces both,
rather than adding a third alongside them.


\begin{table}[h]
\centering
\small
\begin{tabular}{@{}p{3.4cm}p{8.6cm}@{}}
\toprule
\textbf{Formal object} & \textbf{Code locus} \\
\midrule
Admission map $\adm$ (Def.~\ref{def:adm}) & admission-obligation battery evaluated per manufacture request \\
Denial monoid (Con.~\ref{con:denial}) & bit-vector refusal reasons, componentwise \textsc{or} \\
Manufacturing morphism $\muop$ (Def.~\ref{def:mu}) & \code{src/mfg.rs} (RDF $\to$ PDDL) and the Lean~4 autoformalization pipeline (Math $\to$ Lean), the two working $\muop$-pipeline instances documented in \texttt{docs/thesis/math\_manufacturing/} \\
Frame / receipt chain (Def.~\ref{def:receipt}) & BLAKE3 chained receipt, genesis-folded \\
Projection $\Proj$ (Cor.~\ref{cor:noncomp}) & bounded verdict/hash/fitness/reason tuple \\
Cost vector (Def.~\ref{def:cost}) & router \code{CostVector}, lexicographic comparator \\
Status byte (Con.~\ref{con:agent8}) & per-agent 8-bit lane vector, word-parallel sweep \\
\bottomrule
\end{tabular}
\caption{Merged implementation map, keystone objects (both source appendices combined; not duplicated).}
\end{table}


\begin{table}[h]
\centering
\small
\begin{tabular}{@{}p{3.4cm}p{8.6cm}@{}}
\toprule
\textbf{Formal object} & \textbf{Code locus} \\
\midrule
Grounding (Def.~\ref{def:ground}) & \code{bcinr-pddl}\textsuperscript{*} grounding loop, \code{PDDL8\_MAX\_GROUND} bound \\
Forward search (Thm.~\ref{thm:bounded-ground}) & \code{GroundProblem::find\_plan}\textsuperscript{*}, BFS over ground-atom sets \\
Marking / firing (Def.~\ref{def:net}, Prop.~\ref{prop:safe}) & \code{PowlReplayVerifier}\textsuperscript{*} branchless bitset update \\
Farkas certificate (Thm.~\ref{thm:farkas}) & nonconformance witness $\bm y$, $O(p \cdot |T|)$ check \\
Fitness/precision (Def.~\ref{def:fitness}) & Q16.16 population-count ratios, replay accumulators \\
Separable decomposition (Thm.~\ref{thm:sep}) & POWL~2.0 conversion, \code{temporal\_plan\_to\_powl\_tape}\textsuperscript{*} \\
Datalog saturation (App.~\ref{ch:datalog}) & \code{pddl-index} / Nemo reachability index \\
\bottomrule
\end{tabular}
\caption{Merged implementation map, planning/conformance geometry objects.
\textsuperscript{*}These four loci (\code{bcinr-pddl}'s
\code{GroundProblem}, \code{bcinr-powl-receipt}'s
\code{PowlReplayVerifier}, and \code{bcinr-pddl}'s
\code{temporal\_plan\_to\_powl\_tape}) live in the sibling \code{bcinr}
git repository, outside praxis's own \texttt{crates/}/\texttt{src/} trees,
reached only via path dependencies in \texttt{Cargo.toml}; \code{pddl-index}
(Datalog saturation row) is praxis's own in-repo crate, and
\texttt{PDDL8\_MAX\_GROUND} is a constant from the shared
\code{wasm4pm-compat} dependency used by both \code{bcinr-pddl} and
praxis's own \code{pddl-index}, whose own \code{IndexedGroundProblem}
type (distinct from \code{bcinr-pddl}'s \code{GroundProblem}) also
enforces it.}
\end{table}

\chapter{Appendix: The Datalog Index Layer}
\label{ch:datalog}


This appendix reports the Datalog index layer as an implementation-level
addendum, not a full chapter of the main argument: unlike every other
construct in Chapters~\ref{ch:pddl}--\ref{ch:cost}, Datalog saturation and
symbol-dictionary encoding do not recur in \emph{projection\_thesis} and
have no cross-reference into the keystone chapters. It is retained
because the source outline records it as evidence-bound corpus content,
not because it is central to this paper's argument.

\section{Nemo saturation}


\begin{definition}[Datalog saturation]\label{def:saturation}
Let $(\mathcal{D}, \mathcal{Ob})$ be a grounded PDDL problem with initial
atom set $\bm m_0$ and ground action set $T$. The Nemo Datalog program
$\Pi$ has one fact per atom in $\bm m_0$ and one rule per ground action;
the saturation $\mathrm{sat}(\Pi)$ is the least fixpoint
$\mathrm{lfp}(T_\Pi)$ of the immediate-consequence operator; an atom $p$
is reachable iff $p \in \mathrm{sat}(\Pi)$.
\end{definition}
\begin{proposition}[Saturation correctness]\label{prop:satcorrect}
$p \in \mathrm{sat}(\Pi)$ iff there exists a sequence of ground actions
whose add-effects include $p$, i.e.\ $p$ is reachable from $\bm m_0$ in
the Petri-net sense.
\end{proposition}
Both verified in both Lean lanes (\texttt{def:saturation},
\texttt{prop:satcorrect}); the completeness half of
\texttt{prop:satcorrect} cites the standard Datalog fixpoint theorem
(Kowalski 1974, Lloyd 1987) as an external result, not an in-paper
derivation.

\section{Symbol encoding}


\begin{definition}[Symbol dictionary]\label{def:symdict}
A symbol identifier $\mathrm{SymId} \in \mathbb{Z}_{>0}$ is a compact
integer encoding of a ground atom or object name; this paper's own cited
corpus (\texttt{docs/thesis/rdf/03\_planning\_geometry.ttl},
\texttt{math:def\_symdict}, matching \texttt{03\_planning\_geometry.tex})
states that \code{pddl-index} maintains a bidirectional map
\code{SymDict}: $\{\text{strings}\} \leftrightarrow \mathbb{Z}_{>0}$,
encoded forward as \code{FxHashMap<String, u32>} and reverse as
\code{Vec<String>}; zero is reserved as undefined. The live
\texttt{crates/pddl-index/src/dict.rs} disagrees with its own declared
corpus on two points, undisclosed until now: the struct is named
\code{Dict}, not \code{SymDict} (declared \texttt{crates/pddl-index/src/dict.rs:32}),
and its forward map is a \texttt{std} \code{HashMap<String, u32>}
(field \texttt{index}, \texttt{crates/pddl-index/src/dict.rs:36}), not
\code{FxHashMap}. This paper states that discrepancy plainly, in the
same spirit as its disclosures for \texttt{thm:mrr}/\texttt{thm:total}
elsewhere, rather than silently substituting the live-code names for the
corpus's.
\end{definition}
Verified in both Lean lanes (\texttt{def:symdict}).
\begin{proposition}[Dictionary cost]\label{prop:dictcost}
With symbol dictionary encoding, a ground atom $p(o_1,\dots,o_k)$ is a
tuple of $k+1$ \code{SymId} values; equality reduces to integer
comparison, replacing $O(|T| \cdot k \cdot N)$ string comparisons with
$O(1)$-time integer counterparts.
\end{proposition}
Verified in the bare-core lane; recorded \texttt{blocked} (not attempted)
in the Mathlib lane (\texttt{prop:dictcost}).
\begin{construction}[Typed relational join]\label{con:join}
Type-constrained grounding is a relational join: for a schema parameter
$v$ of type $\tau$, the admissible objects are $\{o :
\code{TypeIndex::satisfies}(o, \tau)\}$. This paper's own cited corpus
(\texttt{docs/thesis/rdf/03\_planning\_geometry.ttl}, \texttt{math:con\_join},
matching \texttt{03\_planning\_geometry.tex}) states this as
\code{TypeIndex::satisfies}($\mathrm{SymId}(o), \mathrm{SymId}(\tau)$),
operating on $\mathrm{SymId}$ values. The live method disagrees with that
corpus, undisclosed until now: its real signature is
\code{fn satisfies(\&self, obj: \&str, required: \&str) -> bool}
(\texttt{crates/pddl-index/src/ground.rs:116}), taking plain
\code{\&str} values, not $\mathrm{SymId}$, and walking a string-keyed
parent chain to depth bounded by the height of $\mathcal{T}$, reducing
effective branching from $N^k$ to a product of per-parameter
type-compatible object counts. This paper states that discrepancy
plainly, in the same spirit as its disclosures for
\texttt{thm:mrr}/\texttt{thm:total} elsewhere, rather than silently
substituting the live-code signature for the corpus's.
\end{construction}
Verified in both Lean lanes (\texttt{con:join}).

\chapter{Speculative Extension: Planetary-Scale Projection}
\label{ch:speculative}


Every claim in this chapter is motivated by, but not derived from or
checked against, the evidence-bound chapters above, with one stated
exception detailed below: \texttt{thm:mrr}'s underlying
combinatorial-optimization statement was itself checked, and passed, in
the bare-core Lean lane, and that identical statement already appears,
undisclosed as speculative, in \emph{Mission Geometry}'s own corpus RDF
(\S\ref{ch:speculative} below states this in full -- it is the same
statement under one label, not a coincidental name collision between two
different theorems). Each claim is wrapped in a dedicated
\texttt{speculative} environment so a reader can mechanically distinguish
it from Chapters~\ref{ch:keystone}--\ref{ch:datalog}. None of these
claims should be read at the evidentiary standing of a theorem verified
and placed by its own source in a technical chapter.

\section{Maximum Reachable Revenue: a market-framed extrapolation}


\begin{speculative}[Maximum Reachable Revenue]\label{thm:mrr}
Let $A$ be a set of independent client accounts, $\mathrm{realized}(a)$
the revenue of account $a$ under a target stage, and
$\mathrm{lawful\_targets}(a)$ its evidence-gated target stages. The joint
plan optimization decomposes linearly: $\max_{\text{joint plan}}
\sum_{a \in A} \mathrm{realized}(a) = \sum_{a \in A}
\max_{t \in \mathrm{lawful\_targets}(a)} \mathrm{realized}(a,t)$, reducing
search complexity from $\Omega(\prod_a |T_a|)$ to $O(|A| \cdot |T_a|)$.
\end{speculative}
This is a market/revenue-framed claim layered onto an otherwise
mechanism-level chapter: ``independent client accounts'' and ``realized
revenue'' are a business-scale extrapolation of a combinatorial-
optimization fact (additive separability of a max over independent
coordinates), not something evidenced by \code{praxis} code the way
Chapter~\ref{ch:cost} and earlier are. Its placement here is this
paper's own editorial reclassification on business-framing grounds, not a
hedge \emph{Mission Geometry} itself applies: the source states
\texttt{thm:mrr} as an unhedged theorem, formatted identically to its
neighbors \texttt{thm:lex}/\texttt{cor:dominate}/\texttt{def:makespan},
immediately after \texttt{cor:dominate} in Chapter~\ref{ch:cost}'s own
source material (\texttt{docs/thesis/03\_planning\_geometry.tex}, lines
762--772). Like Theorem~\ref{thm:lang-correct} above, \texttt{thm:mrr}
carries no proof body at all in that source \code{.tex} -- no
\code{begin\{proof\}} block -- per the corpus RDF's own extraction
notes (\texttt{docs/thesis/rdf/03\_planning\_geometry.ttl}, lines
21--26), a gap this paper discloses here on the same standard already
applied to \texttt{thm:lang-correct}. \texttt{thm:mrr} is recorded
\texttt{verified} in the bare-core Lean lane (3 attempts) and
\texttt{unformalized} in the Mathlib lane after a \code{Finset.le\_sup'}
motive-unification failure (3 attempts); this same statement -- not a
differently-scoped label that happens to share a name -- already appears,
undisclosed as speculative, as an ordinary \code{dependsOn}-chained
statement in \emph{Mission Geometry}'s own corpus RDF
(\texttt{docs/thesis/rdf/03\_planning\_geometry.ttl}, tagged
\code{fromPaper "03\_planning\_geometry"} in
\texttt{docs/thesis/rdf/corpus.ttl}), immediately following
\texttt{cor:dominate}/\texttt{thm:lex} and depending on
\texttt{def:cost} -- the same technical cost-arbitration chain
Chapter~\ref{ch:cost} above cites as evidence-bound. This paper's own
declared source corpus therefore does not itself separate this claim from
its technical statement chain, in tension with
\S\ref{sec:thesis-statement}'s promise that speculative material is
never mixed silently into a technical chapter; this paper states that
tension plainly here rather than reading it as a label collision between
two unrelated theorems. A distinct label,
\texttt{thm:mrrsep}, names a separate additive-separation statement and is
recorded the other way around -- \texttt{unformalized} in the bare-core
lane (3 attempts, 4 kernel errors on the final attempt including two
\code{omega} failures) and \texttt{verified} in the Mathlib lane (2
attempts). \texttt{thm:mrrsep} belongs to
\texttt{docs/thesis/rdf/01\_admission\_algebra.ttl}, outside this paper's
own two-source scope (\texttt{03\_planning\_geometry.tex},
\texttt{projection\_thesis.tex}), the same disclosure this paper makes
for \texttt{thm:total} in Chapter~\ref{ch:limitations}; it is mentioned here only
because its label sits adjacent to \texttt{thm:mrr} in the corpus's Lean
receipts, not because it is in scope. These are two distinct Lean files
with two distinct, non-overlapping receipt histories; this paper does not
conflate them under a single ``MRR is verified'' claim.

\section{Planetary arithmetic beyond the exact count}


\begin{speculative}[Planetary arithmetic, estimates]
For a fleet of $N$ agents (Construction~\ref{con:agent8}): compressed
status is $\sim 10\,\mathrm{GB}$ and sweep time is $\Theta(N/(8B))$ at
streaming bandwidth $B$ -- labeled by the source itself as an engineering
estimate, not a theorem. The affordability ratio between one cache-line
mask operation and one LLM admission decision is $\sim 10^{7}$ -- an
order-of-magnitude estimate, explicitly not a theorem. Consequence, not
theorem: a comprehension-based or per-agent-LLM planetary control plane
is infeasible by roughly $N \cdot 10^{7}$ against a feasible bit-parallel
admission sweep; projection-based admission is the only route to
planetary coherence the source argues for at $N = 10^{12}$ agents.
\end{speculative}
The source paper's own labeling -- ``engineering estimate,''
``order-of-magnitude estimate,'' ``consequence, not theorem'' -- is kept
verbatim here rather than upgraded to theorem language. The receipts for
these sub-claims (\texttt{prop:planetary} parts (b)--(d)) are recorded
\texttt{verified} in both Lean lanes, but that status attaches to the
proposition \emph{wrapper} the source paper itself states parts (b)/(c)
are non-theorem estimates within -- a Lean receipt on the wrapper does not
convert an estimate into a machine-checked numeric claim, and this paper
does not read it as doing so.

\section{A shared refusal vocabulary across agents}


\begin{speculative}[A refusal vocabulary legible across agents]
Refusal categories, drawn from Construction~\ref{con:denial}'s $n$-lane
denial monoid and organized under the source's fixed eight-bucket
surjection, are claimed, in the source's own words, to form a vocabulary
such that ``a refusal category is legible to a recipient who never saw
the sender's interior'' at trillion-agent scale.
\end{speculative}
This is an interpretive/legibility claim about a fixed enumeration of
refusal categories, not a corollary Construction~\ref{con:denial} or
Proposition~\ref{prop:semilattice} proves; those establish the algebraic
structure of denial words, not that any particular 8-bucket enumeration is
legible across arbitrary agents. It is recorded \texttt{verified} in both
Lean lanes under \texttt{con:denial}'s own label, but that receipt
attaches to the monoid construction, not to the legibility claim layered
on top of it, which is why this paper states the legibility claim only
here, quarantined, rather than in Chapter~\ref{ch:keystone}.

\section{The self-audited-scale antibody clause}


\begin{speculative}[Self-audited-scale antibody clause]\label{prop:selfauditedscale}
Let a system emit claims $c$ with receipts $r(c)$, and define the
overclaim set $\mathrm{Over} = \{c : \nexists\, r(c) \text{ with }
V(\Proj) = 1\}$. If the system enforces $\mathrm{Over} = \varnothing$
(the \textsf{docs-exceed-mechanism} defect class), the admissible
magnitude of a claim is bounded by the mechanism that receipts its
limits, not by rhetoric.
\end{speculative}
Recorded \texttt{blocked} with zero attempts in both Lean lanes
(\texttt{prop:selfauditedscale}). It is stated here as a design aspiration
this paper itself is trying to hold to -- every claim above names its
artifact or is marked speculative -- not as a proved property of any
running system.

\chapter{Limitations}
\label{ch:limitations}


This chapter states what is not shown, following \texttt{CLAUDE.md}'s
reporting instruction to state findings, not verdicts, and without the
quality adjectives the task does not ask for.


Three central keystone theorems -- \texttt{thm:faithful},
\texttt{cor:noncomp}, \texttt{thm:conservation} -- are stated in
Chapter~\ref{ch:keystone} but are not machine-verified in either Lean
lane on the receipts currently on disk: \texttt{thm:faithful} has no
bare-core receipt line and is recorded \texttt{excluded} in the Mathlib
lane; the other two are recorded \texttt{blocked} with zero attempts in
both lanes. Chapter~\ref{ch:instance}'s trillion-agent instance argument
depends on \texttt{thm:faithful} and inherits this exact gap.


Farkas separation (\texttt{thm:farkas}) remains genuinely open in the
Mathlib lane: the missing lemma is named precisely
(\code{PointedCone.FG.isClosed}, the topological half of
Minkowski--Weyl), not a vague timeout, and is recorded
\texttt{blocked}/not-attempted in the bare-core lane entirely. Two related
Petri-net propositions, \texttt{prop:cone} (bare-core) and
\texttt{thm:bounded-ground} (Mathlib), each regressed against the other
lane for reasons specific to that lane's own Mathlib-snapshot lookup
failures (Ch.~\ref{ch:pddl}, Ch.~\ref{ch:petri} above) -- these are
lane-specific library gaps, not defects shared across both formalizations
of the same statement.


The Mathlib-linked lane covers 178 of the corpus's 202 distinct labels
(88\%), not all of it, and not a claim of parity with the
separately-verified bare-core lane's own scope; the two lanes are
independently verified over overlapping but non-identical statement sets
(\texttt{docs/thesis/math\_manufacturing/rdf/thesis.ttl}, Ch.~4, reports
this same 178/202 figure for the full corpus, of which this paper's own
planning-geometry and projection-keystone labels are a subset). Three
further Mathlib-lane failures exist corpus-wide beyond
\texttt{thm:farkas} and \texttt{thm:bounded-ground} and fall within this
document's own two-source scope: \texttt{def:attnbalance},
\texttt{prop:conformembership}, and \texttt{thm:mrr} -- each has a
specific, named kernel error on file (a library-lookup or motive-unification
gap), not a generic failure; none is silently retried past the fixed
3-attempt bound. A fourth label, \texttt{thm:total}, is reported
\texttt{unformalized} in the same Mathlib-migration receipts file but is
not the same kind of gap: its recorded \code{last\_error}
(\texttt{mathlib\_migration\_receipts.jsonl:199}) states plainly that the
theorem's own claim -- \code{Category.Reserved} has empty preimage under
\code{RefusalScenario.category} -- is false against the already-verified
\texttt{def:tax}, which maps three \code{RefusalScenario} variants
(\code{logicContradiction}, \code{andonPull}, \code{andonEscalation}) to
\code{Category.Reserved}
(\texttt{Praxis/Corpus/def\_tax.lean:67--69},
\texttt{Praxis/Corpus/thm\_total.lean:19}); no amount of tactic work
closes that goal because the goal itself is false against a dependency
already on the kernel's own record, not merely difficult. This is a
demonstrated falsification, not a tooling gap, and \texttt{thm:total}
belongs to \texttt{docs/thesis/rdf/01\_admission\_algebra.ttl}, outside
this paper's own two-source scope (\texttt{03\_planning\_geometry.tex},
\texttt{projection\_thesis.tex}). Unresolved and stated here rather than
smoothed over: the bare-core lane
(\texttt{tools/paper-factory/lean-pilot/formalization\_receipts.jsonl:192})
records this same label \texttt{thm:total} as \texttt{verified} (3
attempts, no error) against the identical falsified claim -- an
unremarked cross-lane discrepancy of the same kind this paper flags
elsewhere for \texttt{thm:mrr}/\texttt{thm:mrrsep}
(Chapter~\ref{ch:speculative}) and \texttt{thm:farkas}/
\texttt{prop:conformembership} (Chapter~\ref{ch:conformance}); this paper
does not resolve that discrepancy, only reports it.


This is a single-kernel result: only Lean~4's own kernel accepted or
rejected each statement; no independent proof assistant cross-checks any
of the \texttt{verified} dispositions reported in this paper.
\texttt{thm:lang-correct} (Ch.~\ref{ch:separability}) is flagged
separately: its source \texttt{.tex} carries no proof body, so its status
here is asserted-with-a-Lean-formalization-on-record, a weaker claim than
proved-in-paper, stated as such rather than smoothed into an unqualified
``proved.''


The specific rhetorical \emph{device} of \emph{projection\_thesis}'s own
front matter -- addressed to $10^{12}$ concurrent agents, a
bounded-context-capacity motif citing Miller and Cowan -- is deliberately
not carried forward as technical content anywhere in this paper outside
the explicit disclaimer of Chapter~\ref{ch:speculative}: the source
outline itself marks it speculative and instructs against duplicating it
as if it were a technical claim, and this paper follows that instruction
rather than restating the numeric/motif device a third time across the
thesis series. This is narrower than the chapter's own \emph{naming}
practice: Chapter~\ref{ch:instance}'s title
(``Instance: An Autonomous Build and the Trillion-Agent Status Byte'')
and this paper's own introductory structure overview (Chapter~1) and the
paragraph immediately above in this same chapter all use ``trillion-agent''
as the name of that chapter's status-byte instance, not as a restatement
of the $10^{12}$/Miller--Cowan motif; that naming usage is excepted from
this paragraph's
claim and is not itself disclosed as speculative, since it names an
evidence-bound instance argument (Chapter~\ref{ch:instance}), not the
front-matter framing device.

\chapter{Conclusion}
\label{ch:conclusion}


This paper states one evidence-bound keystone (Chapter~\ref{ch:keystone})
and its concrete planning-and-conformance geometry
(Chapters~\ref{ch:pddl}--\ref{ch:cost}), one evidence-bound instance at
declared scope (Chapter~\ref{ch:instance}), and a clearly quarantined set
of estimates and extrapolations (Chapter~\ref{ch:speculative}) that this
paper does not claim are anything more than what their own source labels
them. The correct adversarial question after this consolidation is not
whether gaps exist -- Chapter~\ref{ch:limitations} names them plainly --
but which specific admitted statement, blocked dependency, or
evidence-bound claim above is wrong, checkable against the two receipts
files and the corpus RDF that produced it, both cited by exact path
throughout.


